% SPDX-License-Identifier: GPL-3.0-or-later
% Copyright (C) 2026 Alessandro Greco
\chapter{Arti: Tor на Rust}\label{cap:30-arti}
\aggiornatoal{23 сентября 2026 г.}{Arti 2.6.0 (1 сентября 2026 г.); tor 0.4.9.13; лабораторная сеть}

\begin{obiettivi}
К концу этой главы вы сможете:
\begin{itemize}[nosep]
  \item объяснять, почему Tor Project переписывает tor на Rust, и оценивать прогресс проекта;
  \item понимать архитектуру Arti и находить в ней концепции, рассмотренные в Части~II;
  \item интерпретировать таблицу поддерживаемых подпротоколов и понимать, какие функции всё ещё отсутствуют;
  \item настраивать Arti как клиент приватной сети книги (Лабораторная работа~20).
\end{itemize}
\end{obiettivi}

На протяжении всей книги мы опирались на \emph{tor}, реализацию на C, разработка которой началась в 2002 году. Однако будущее клиентской части Tor — \textbf{Arti}, полная реализация с нуля на Rust. Tor Project обозначил Arti как предпочтительный клиент для разработки новой функциональности и объявил, что каждая новая клиентская функция будет получать реализацию в Arti. Эта глава рассматривает Arti глазами читателя, который в лабораторных работах только что написал небольшой Tor-клиент на Python. Многие архитектурные решения Arti — те же решения, с которыми столкнулись мы, только реализованные в промышленном масштабе.

% ---------------------------------------------------------------------------
\section{Зачем переписывать tor?}\label{sec:arti-perche}

Вводный документ проекта перечисляет четыре причины:
\begin{description}
  \item[Безопасность] Безопасный код Rust предотвращает многие ошибки работы с памятью, но небезопасный код, зависимости и логика приложения всё ещё требуют проверки. По оценке Tor Project, как минимум половина зарегистрированных уязвимостей безопасности tor была бы невозможна на Rust, а многие другие стали бы крайне маловероятными.
  \item[Скорость разработки] Гарантии на уровне языка позволяют разрабатывать новый код увереннее и быстрее.
  \item[Гибкость] tor начинался как прокси SOCKS, а возможности интеграции добавлялись позднее. Arti изначально проектировался как модульная библиотека, которую можно встраивать в другие приложения.
  \item[Сопровождаемость] Двадцать лет разработки создали сложные взаимосвязи между компонентами кодовой базы tor. Переписывание даёт возможность яснее разделить ответственности.
\end{description}

\paragraph{Вехи.} Первая публичная версия Arti, 0.0.1, появилась в октябре 2021 года. Версия 1.0.0, выпущенная 1 сентября 2022 года, объявила Arti стабильным клиентом. В августе 2023 года версия 1.1.7 включила сборку клиентской поддержки onion-сервисов по умолчанию, хотя пользователю всё ещё требовалось явно активировать функцию. Версия 2.0.0 вышла в феврале 2026 года. Эталонная версия этой главы — 2.6.0 от 1 сентября 2026 года; в ней управление перегрузкой (\cref{sec:cc}) и шифрование CGO (\cref{sec:cgo}) включены постоянно.

\paragraph{Текущая область применения.} Согласно документации проекта на эталонную дату, Arti — полноценный универсальный клиент с поддержкой мостов и подключаемых транспортов; он умеет и обращаться к onion-сервисам, и публиковать их. Он пока не может работать как ретранслятор, мост или центр каталогов. Эти возможности находятся в разработке, и значительная часть работы в последних релизах была сосредоточена на поддержке ретрансляторов и центров каталогов.

% ---------------------------------------------------------------------------
\section{Архитектура}\label{sec:arti-architettura}

Arti состоит из десятков \emph{крейтов} Rust — единиц компиляции языка: в версии 2.6.0 их более семидесяти. Соглашения об именовании различают роли. Пользовательские крейты начинаются с \texttt{arti}, специфичные для протокола Tor — с \texttt{tor-}, а крейты, полезные вне Tor, имеют более общие названия. Архитектуру определяет важный принцип: код, который напрямую читает или записывает сетевые данные, должен быть минимальным; большинство компонентов должно вести себя \enquote{так, как если бы сети не существовало}, что упрощает тестирование.

\Cref{fig:crate} сопоставляет основные крейты с модулями \texttt{torlab}, выполняющими сходные задачи. Соответствие не случайно: обе архитектуры следуют уровням протокола, от низкоуровневых примитивов до высокоуровневого клиента.

\begin{figure}[htbp]
\centering
% SPDX-License-Identifier: GPL-3.0-or-later
% Copyright (C) 2026 Alessandro Greco
%% Русское альтернативное описание EPUB.
\altfig{Учебная схема некоторых крейтов Arti: сверху программа arti с SOCKS-прокси и arti-client, посередине управление каталогами, цепями, сторожевыми узлами и onion-сервисами, снизу протокол, ячейки и криптография. Сплошные стрелки показывают упрощённые функциональные связи, а не точный граф зависимостей Cargo. Пунктирные линии сравнивают задачи с учебными модулями torlab (directory.py, path\_selection.py, onion.py, channel.py, circuit.py, cells.py, ntor.py, encryption.py), не обозначая зависимостей или полной эквивалентности.}
% Основные crates Arti и соответствующие модули torlab
\begin{tikzpicture}[font=\sffamily\scriptsize,
  crate/.style={draw=torviola,fill=torviola!8,rounded corners=2pt,minimum width=2.5cm,minimum height=0.6cm},
  lab/.style={draw=torverde!60!black,fill=torverde!10,rounded corners=2pt,minimum width=2.2cm,minimum height=0.6cm,font=\ttfamily\scriptsize}]
  \node[crate] (arti) at (0,4.2) {arti (SOCKS, CLI)};
  \node[crate] (client) at (0,3.3) {arti-client};
  \node[crate] (dirmgr) at (-2.9,2.3) {tor-dirmgr};
  \node[crate] (circmgr) at (0,2.3) {tor-circmgr};
  \node[crate] (guardmgr) at (2.9,2.3) {tor-guardmgr};
  \node[crate] (netdoc) at (-2.9,1.3) {tor-netdoc};
  \node[crate] (chanmgr) at (0,1.3) {tor-chanmgr};
  \node[crate] (hs) at (2.9,1.3) {tor-hsclient};
  \node[crate] (proto) at (0,0.3) {tor-proto};
  \node[crate] (cell) at (0,-0.7) {tor-cell};
  \node[crate] (llc) at (0,-1.7) {tor-llcrypto};
  \draw[->] (arti) -- (client);
  \draw[->] (client) -- (dirmgr); \draw[->] (client) -- (circmgr); \draw[->] (client) -- (hs);
  \draw[->] (dirmgr) -- (netdoc); \draw[->] (circmgr) -- (chanmgr); \draw[->] (circmgr) -- (guardmgr);
  \draw[->] (chanmgr) -- (proto); \draw[->] (proto) -- (cell); \draw[->] (cell) -- (llc);
  \draw[->] (dirmgr) -- (circmgr); \draw[->] (hs) -- (circmgr);
  % Концептуальная схема: пунктир — аналогии, а не зависимости Cargo.
  \node[lab] (d) at (-6.2,1.3) {directory.py};
  \node[lab] (p) at (6.0,2.3) {path\_selection.py};
  \node[lab] (o) at (6.0,1.3) {onion.py};
  \node[lab] (c) at (-4.4,0.3) {channel.py, circuit.py};
  \node[lab] (ce) at (-4.4,-0.7) {cells.py};
  \node[lab] (n) at (-4.4,-1.7) {ntor.py, encryption.py};
  \foreach \a/\b in {d/netdoc,p/guardmgr,o/hs,c/proto,ce/cell,n/llc} \draw[dashed,torgrigio] (\a) -- (\b);
  \node[torviola] at (0,5) {Arti (Rust)};
  \node[torverde!50!black] at (-5.2,2.4) {torlab (Python)};
\end{tikzpicture}

\caption{Упрощённая функциональная схема крейтов Arti и учебных модулей \texttt{torlab}. Сплошные стрелки показывают архитектурные связи, но не точный граф зависимостей Cargo; пунктир указывает на сходные задачи, а не на зависимости между проектами или полную эквивалентность. Например, \texttt{onion.py} только расшифровывает дескрипторы и не заменяет клиент onion-сервисов; крейты \texttt{tor-bytes} и \texttt{tor-hsservice} опущены.}
\label{fig:crate}
\end{figure}

\begin{description}
  \item[\texttt{tor-llcrypto}, \texttt{tor-bytes}] Криптографические примитивы и разбор/сериализацию двоичных данных. Они предоставляют основу, которую наши \file{ntor.py} и \file{encryption.py} получали из модулей Python \texttt{cryptography} и \texttt{struct}.
  \item[\texttt{tor-cell}] Ячейки и ретрансляционные сообщения: аналог \file{cells.py} в Arti (\cref{cap:08-celle}).
  \item[\texttt{tor-proto}] Каналы, цепи, послойное шифрование и потоки: всё, что мы строили с помощью \file{channel.py}, \file{circuit.py} и \file{stream.py}.
  \item[\texttt{tor-chanmgr}, \texttt{tor-circmgr}, \texttt{tor-guardmgr}] Управление множеством каналов и цепей, выбор маршрутов и поддержка сторожевых узлов. Наш учебный клиент почти не реализовывал эти возможности, хотя они относятся к наиболее сложным частям настоящего клиента.
  \item[\texttt{tor-netdoc}, \texttt{tor-dirmgr}] Чтение и обновление документов каталога, соответствующие \file{directory.py}.
  \item[\texttt{tor-hsclient}, \texttt{tor-hsservice}] Обе стороны onion-сервисов; наш \file{onion.py} реализовывал только расшифрование дескриптора.
  \item[\texttt{arti-client}, \texttt{arti}] Высокоуровневая библиотека для встраивания в другие приложения и командная программа, предоставляющая прокси SOCKS.
\end{description}

\subsection{Сравнение: запись ячейки}

В Лабораторной работе~2 мы написали функцию кодирования ячейки. Соответствующая функция Arti, \texttt{write\_cell} в \file{crates/tor-cell/src/chancell/codec.rs}, выполняет те же шаги: записывает CircID, затем команду. Для ячейки переменной длины резервирует два байта под длину, кодирует тело, затем возвращается и записывает фактическую длину. Для ячейки фиксированной длины кодирует тело и проверяет, что оно не превышает 509 байт (\texttt{CELL\_DATA\_LEN}). Два отличия показательны:
\begin{itemize}
  \item CircID \emph{всегда} записывается четырьмя байтами. Arti не реализует версии протокола соединения с номером ниже 4, поэтому специальный случай двухбайтовых CircID ему не нужен (\cref{cap:08-celle}). Только ячейка \cella{VERSIONS}, предшествующая согласованию, обрабатывается отдельно.
  \item Нарушение ограничения длины приводит не к ошибке вызывающего кода, а к \emph{внутренней} ошибке: слишком большая исходящая ячейка была бы дефектом самого Arti, а не ошибкой, вызванной данными, полученными от удалённого узла. Тип ошибки явно отличает баг реализации от недоверенного ввода.
\end{itemize}
Arti сознательно не поддерживает устаревшие протоколы. Его цель — совместимость с \emph{рекомендуемыми сейчас} версиями, а не сохранение каждого исторического формата.

% ---------------------------------------------------------------------------
\section{Поддерживаемые функции}\label{sec:arti-supporto}

Самое точное описание возможностей реализации Tor — поддерживаемые ею \textbf{подпротоколы} (\spec{tor-spec}{Subprotocol-based versioning}). Каждый ретранслятор и клиент объявляет поддерживаемые версии для каждой группы функций: \texttt{Link}, \texttt{Relay}, \texttt{FlowCtrl}, \texttt{HSIntro} и т.д. Документация совместимости Arti перечисляет, среди прочего, следующие клиентские возможности:
\begin{itemize}
  \item \texttt{Link=4} и \texttt{Link=5}: четырёхбайтовые CircID и современное заполнение канала (\cref{cap:07-canali});
  \item \texttt{Relay=4} (ntor-v3, \cref{cap:09-circuiti}), \texttt{Relay=5} (согласование подпротоколов) и \texttt{Relay=6} (шифрование CGO);
  \item \texttt{FlowCtrl=1} (аутентифицированные SENDME) и \texttt{FlowCtrl=2} (управление перегрузкой);
  \item \texttt{Desc=4}: семейства на основе ключей, описанные в \cref{sec:relay-famiglie}.
\end{itemize}

\begin{attenzione}[Ошибка на единицу в таблице Arti]
В указанном документе совместимости Arti эти два подпротокола управления потоком пронумерованы как \texttt{FlowCtrl=0} и \texttt{FlowCtrl=1}. Однако спецификация назначает аутентифицированным SENDME значение 1, а управление перегрузкой — 2~\spec{tor-spec}{Subprotocol-based versioning}; ретрансляторы в публичном консенсусе объявляют \texttt{FlowCtrl=1-2}. Это ошибка на единицу в документации, а не в реализации. Она также показывает полезную практику: когда источники расходятся, сверяйтесь со спецификацией и ищите независимое подтверждение в консенсусе.
\end{attenzione}

На дату, принятую в главе за точку отсчёта в клиенте Arti отсутствуют две возможности C-клиента. \texttt{Conflux=1} обеспечивает многопутевую передачу: crate \texttt{tor-proto} поддерживает её, но Arti пока не строит цепи Conflux. \texttt{Padding=2} обозначает автоматы заполнения, предназначенные для маскировки построения цепей onion-сервисов (\cref{sec:att-wf}). Хотя \texttt{tor-proto} содержит механизм заполнения на основе библиотеки \emph{maybenot}, сами машины ещё не реализованы. Ни один подпротокол со стороны ретранслятора пока не помечен поддерживаемым: значительная часть кода для ретрансляторов уже существует, но функции будут объявлены поддерживаемыми только после появления работающего ретранслятора, который способен их предоставлять.

% ---------------------------------------------------------------------------
\section{Лабораторная работа: Arti в приватной сети}\label{sec:lab-arti}

\begin{laboratorio}[20]{Второй клиент — на Rust}
\textbf{Цель.} Скомпилировать Arti, настроить его как клиент приватной сети книги и использовать его прокси SOCKS вместо \texttt{test008c}.\\
\textbf{Предварительные условия.} Лабораторная сеть должна быть запущена; требуется инструментарий Rust (\cmd{cargo}).\\
\textbf{Оценочное время.} 1--2 часа, большая часть времени уходит на компиляцию.
\end{laboratorio}

\paragraph{Шаг 1: компиляция.} Из корня репозитория книги получите эталонную версию Arti 2.6.0 в \file{labs/ru/vendor/arti} и соберите исполняемый файл \texttt{arti}:
\begin{lstlisting}[language=bash,numbers=none]
$ mkdir -p labs/ru/vendor
$ git clone https://gitlab.torproject.org/tpo/core/arti.git labs/ru/vendor/arti
$ git -C labs/ru/vendor/arti checkout arti-v2.6.0
$ (cd labs/ru/vendor/arti && cargo build --release -p arti --locked)
$ labs/ru/vendor/arti/target/release/arti --version
Arti 2.6.0
\end{lstlisting}
Первая сборка занимает несколько минут, поскольку компилирует всё дерево крейтов и их зависимостей.

\paragraph{Шаг 2: конфигурация.} Arti не читает \file{torrc}; он использует TOML. Для доступа к непубличной Tor-сети Arti должен знать центры каталогов и резервные каталожные кеши, из которых получит начальные документы каталога, подобно строкам \texttt{DirAuthority} в \file{torrc}, созданных Chutney (\cref{sec:rp-file}). Скрипт \file{labs/ru/лаборатория/arti\_config.py} читает отпечаток RSA, ключ Ed25519, идентификатор сертификата v3 и порты каждой авторитетной ноды из локальных файлов. Он записывает значения в \file{labs/ru/лаборатория/net/arti/arti.toml}. Следующий сокращённый фрагмент демонстрационный, а не готовый TOML: идентификаторы генерируются локально, а многоточия обозначают пропущенные строки:
\begin{lstlisting}[numbers=none]
[tor_network]
fallback_caches = [
  { rsa_identity = "7B465200C885ED566C6CD4F46DBB0492076D203A",
    ed_identity = "JwhsA9VS01mRVnQseL9xsKZfstoBae7OcgL3jBvDzis",
    orports = ["127.0.0.1:5100"] },
  ...
]
[tor_network.authorities]
v3idents = ["390998046DEFF1E68AC346D7577319791B32C91A", ...]
\end{lstlisting}
Ещё три изменения конфигурации адаптируют Arti к сети, где всё работает на \texttt{127.0.0.1}. \texttt{allow\_local\_addrs} разрешает соединения с локальными адресами, которые клиенты публичной сети обычно отклоняют. Две невозможные длины сетевых префиксов — 33 бита для IPv4 и 129 для IPv6 — отключают ограничение совместного сетевого префикса, иначе блокирующий здесь каждую цепь. Наконец, \texttt{hsdir\_interval} сокращает период onion-сервиса так же, как \opt{TestingTorNetwork} в tor. Эти параметры повторяют конфигурации Arti, генерируемые самим Chutney.

\paragraph{Шаг 3: запуск.} После создания конфигурации командой \file{python3 labs/ru/лаборатория/arti\_config.py} из корня книги выполните \file{labs/ru/vendor/arti/target/release/arti proxy -p 9150 -c labs/ru/лаборатория/net/arti/arti.toml}. Время начальной загрузки зависит от машины; в одном из наших запусков:
\begin{lstlisting}[numbers=none]
INFO tor_dirmgr::bootstrap: 1: Looking for a consensus.
INFO tor_dirmgr::bootstrap: 1: Downloading certificates for consensus
     (we are missing 3/3).
INFO tor_dirmgr::bootstrap: 1: Downloading microdescriptors (we are missing 8).
INFO tor_dirmgr: Marked consensus usable.
WARN tor_netdir: Requested 60 relays, but only 8 were usable, ...  role=Guard
INFO arti::subcommands::proxy: Sufficiently bootstrapped; proxy now functional.
INFO tor_guardmgr::guard: We have found that guard [scrubbed] is usable.
\end{lstlisting}
Три детали показательны. Процедура начальной загрузки следует \cref{sec:dir-client} буквально по шагам: сначала консенсус, затем сертификаты трёх центров каталогов и наконец восемь микродескрипторов. Предупреждение \texttt{tor\_netdir} отражает размер нашей крошечной сети: менеджер сторожевых узлов предпочёл бы выборку из 60 кандидатов, но доступны только восемь. Идентификатор сторожевого узла в этом примере заменён на \texttt{[scrubbed]}: Arti помечает чувствительные значения с помощью крейта \texttt{safelog} и скрывает их в логах, аналогично \opt{SafeLogging} в tor. В других журналах обновлений может появляться \texttt{Applying a consensus diff} (в приведённом отрывке этой строки нет). При последующих обновлениях Arti может скачать только изменения, если доступна разница, согласованная с помощью \texttt{DirCache=2}; такой ответ не гарантируется при каждом обновлении.

\paragraph{Шаг 4: использование.} Arti предоставляет прокси SOCKS на порту 9150. \cmd{curl} использует его так же, как раньше \texttt{test008c}:
\begin{lstlisting}[language=bash,numbers=none]
$ curl -s -o /dev/null -w '%{http_code} in %{time_total}s\n' \
       --noproxy "" --socks5-hostname 127.0.0.1:9150 http://127.0.0.1:8080/
200 in 0.047224s
\end{lstlisting}
Пустой аргумент \texttt{-{}-noproxy} не позволяет переменной \texttt{NO\_PROXY} обходить SOCKS для локального адреса. В предполагаемом тесте цепь проходит через ретрансляторы на C в приватной сети: клиент Rust и ретранслятор на C говорят на одном и том же байтовом протоколе, изученном в Части~II. Последнее наблюдение связано с Лабораторной работой~18: SOCKS-запросы, содержащие только IP-адрес (\texttt{-{}-socks5}), также успешно выполняются, а стандартные журналы Arti не показывают предупреждение, аналогичное tor. Однако документация крейта \texttt{arti-client} предупреждает пользователей API, что локальное разрешение имён может раскрыть назначение DNS-серверу. Сам по себе SOCKS-запрос с IP-адресом не доказывает утечку DNS (\cref{sec:lab-opsec}).

Лабораторная работа входит в \texttt{make verifica-lab}, если Arti собран; иначе проверка пропускается с соответствующим сообщением.

% ---------------------------------------------------------------------------
\section*{Итоги}
\begin{itemize}
  \item Arti переписывает tor на Rust с механизмами безопасности памяти. Он модульный, проектировался как библиотека, стабилен как клиент с 2022 года и является предпочтительным клиентом для разработки новых функций.
  \item Его архитектура следует уровням протокола: примитивы и ячейки, каналы и цепи, затем управление каталогом, маршрутами и guard.
  \item Версии подпротоколов описывают возможности Arti. Клиент реализует ntor-v3, CGO, управление перегрузкой и семейства узлов на основе ключей, тогда как Conflux и автоматы заполнения на эталонную дату главы ещё не завершены. Поддержка роли ретранслятора всё ещё разрабатывается.
  \item Arti подключается к приватной сети через настроенные центры каталогов и резервные кеши, подобно узлу Chutney.
\end{itemize}

% ---------------------------------------------------------------------------
\section*{Упражнения}

\begin{esercizio}[1]\label{es:30-crate}
Для каждой главы Части~II определите крейт Arti, реализующий основные рассматриваемые там понятия.
\end{esercizio}

\begin{esercizio}[1]\label{es:30-obsoleto}
Каковы преимущества и недостатки решения Arti не реализовывать устаревшие версии протокола? Кто может оказаться исключён?
\end{esercizio}

\begin{esercizio}[2]\label{es:30-cella}
Прочитайте \texttt{decode\_cell} в \file{crates/tor-cell/src/chancell/codec.rs} и сравните с \texttt{decode} из Лабораторной работы~2. Как распознаётся ячейка переменной длины и что происходит, если входных данных ещё недостаточно для полной ячейки?
\end{esercizio}

\begin{esercizio}[2]\label{es:30-isolamento}
Повторите эксперимент по изоляции потоков из Лабораторной работы~11, используя Arti как прокси SOCKS. Не позволяют ли разные учётные данные SOCKS объединять потоки в одну цепь в Arti? Отличайте изоляцию от ошибочного предположения, что каждой паре данных навсегда соответствует ровно одна цепь; найдите параметры настройки.
\end{esercizio}

\begin{esercizio}[3]\label{es:30-libreria}
Напишите небольшую программу на Rust с \texttt{arti-client}, которая получает страницу с сервера лабораторной сети, загружая конфигурацию Лабораторной работы~20 из файла.
\end{esercizio}

% ---------------------------------------------------------------------------
\section*{Дополнительная литература}
\begin{itemize}
  \item Сайт проекта \url{https://arti.torproject.org/} и репозиторий исходного кода, особенно README, \emph{обзор архитектуры} (\file{doc/dev/Architecture.md}) и руководство совместимости (\file{doc/Compatibility.md}).
  \item Журнал изменений Arti, документирующий изменения каждого ежемесячного релиза.
  \item Конфигурация Arti, генерируемая Chutney (\file{lib/chutney/arti/config.py}), на которой основана Лабораторная работа~20.
\end{itemize}
